\section{Implementation and Reproducibility}
\label{sec:impl}

\paragraph{Code organization}
The experiment pipeline is implemented under the project source directory with
modular components for configuration, data loading, and autoregressive baselines,
speculative decoding, metric aggregation, and runtime orchestration.
The notebook driver executes these modules in fixed phase order and writes
artifacts to the results directory.

\paragraph{Model loading and precision}
In the current configuration, target and draft models are loaded in 16-bit
floating point (FP16) for
the main sweep, with quantization controls retained in configuration for
alternative runs. The tokenizer and generation limits are kept constant between
baselines and speculative runs to preserve paired comparability.

\paragraph{Baselines}
We run one autoregressive baseline per regime (deterministic and stochastic)
with the same manifests and output-length constraints as speculative runs.
Each baseline artifact stores per-sample latency, output length, and task
predictions.

\paragraph{Speculative decoding implementation.}
The speculative loop follows standard token-level draft--verify logic.
At each cycle, the draft proposes up to $k$ tokens autoregressively; the
target verifies the proposal in one pass over the extended prefix; the accepted
prefix is committed up to the first mismatch; and one correction token is added
from the target distribution when needed. This procedure preserves target-level
generation semantics while exposing acceptance statistics needed for analysis.
Both draft and target states are tracked with key-value (KV) caches throughout
the decode loop.
Figure~\ref{fig:specdec-impl-flow} summarizes this implementation flow.

\paragraph{CUDA-graph runtime note}
Graph-capture metadata differs by run family. The Qwen3 RTX4090 summary
records captured execution with capture rate 1.0, whereas the Qwen2.5
RTX4090 summary records pending capture with rate 0.0. RTX5090-laptop
sensitivity runs do not show a consistent latency benefit from graph-enabled
toggles. CUDA graphs are therefore reported as runtime context, not as an
independently evaluated optimization.

\paragraph{Dynamic control flow}
Acceptance-dependent branches and cache updates require host-visible
orchestration in the native PyTorch loop and can prevent full graph capture.
Because this overhead was not isolated experimentally, we treat it as a
possible explanation for runtime sensitivity rather than a measured cause.

\paragraph{Tokenizer scope}
Same-family pairings use matching tokenizers and avoid vocabulary alignment
inside verification. Cross-tokenizer pairings are outside this study; their
mapping overhead and acceptance behavior require separate evaluation.

\begin{figure}[t]
\centering
\resizebox{0.98\columnwidth}{!}{%
\begin{tikzpicture}[
	node distance=4.8mm and 6mm,
	box/.style={rectangle, draw, rounded corners=1mm, align=center, minimum height=6mm, text width=31mm, inner sep=1.5mm},
	decision/.style={diamond, draw, align=center, aspect=2, inner sep=1.1mm, text width=22mm},
	line/.style={-Latex, thick, shorten >=1pt, shorten <=1pt}
]
\node[box] (start) {Start sample\\encode prompt and init KV cache};
\node[box, below=of start] (draft) {Draft model proposes\\up to $k$ tokens};
\node[box, below=of draft] (verify) {Target verifies draft block\\in one forward pass};
\node[decision, below=of verify] (accept) {All $k$ tokens\\accepted?};
\node[box, left=of accept, xshift=-9mm] (bonus) {Commit $k$ tokens\\+ one target bonus token};
\node[box, right=of accept, xshift=9mm] (correct) {Commit accepted prefix\\+ one corrected target token};
\node[box, below=of accept, yshift=-2mm] (crop) {Crop/update draft and target\\KV caches to committed prefix};
\node[decision, below=of crop] (stop) {EOS or max tokens\\reached?};
\node[box, below=of stop] (end) {Finish sample and\\record runtime/quality fields};

\draw[line] (start) -- (draft);
\draw[line] (draft) -- (verify);
\draw[line] (verify) -- (accept);
\draw[line] (accept) -- node[above, font=\scriptsize] {yes} (bonus);
\draw[line] (accept) -- node[above, font=\scriptsize] {no} (correct);
\draw[line] (bonus) |- (crop.west);
\draw[line] (correct) |- (crop.east);
\draw[line] (crop) -- (stop);
\draw[line] (stop) -- node[right, font=\scriptsize] {yes} (end);
\draw[line] (stop.west) -- ++(-14mm,0) |- node[pos=0.25, left, font=\scriptsize] {no} (draft.west);
\end{tikzpicture}%
}
\caption{Vanilla speculative decoding implementation flow.
The target model remains the verifier of record; accepted draft prefixes are
committed, and caches are cropped to preserve exact target-consistent decoding.}
\label{fig:specdec-impl-flow}
\end{figure}

\paragraph{Regimes and outputs}
The same verification core is used for deterministic and stochastic modes,
with only the sampling policy changed by regime parameters.
Each configuration writes a dedicated CSV artifact with per-sample runtime,
token counts, and quality-relevant outputs.
